\section{Lean formalization plan}\label{sec:lean}

The archive contains \code{Lean/BringQuinticBlueprint.lean} and a separate
roadmap.  The file is intentionally a blueprint rather than a claim of a new
completed build.  This section gives the precise path from the human proofs to
closed terms of the same form as the existing ProveIt derivations.

\subsection{Target theorem shape}

The existing derivation library represents a closed \ipra{} theorem as
\[
 ([]:\operatorname{List}(\mathrm{Fm}\;\mathrm{PTm}\;0))
 \;\vdash[\mathrm{PRAAx}]\;\varphi.
\]
A new theorem should have this type after all free variables are universally
closed.  The metatheoretic Lean proof may use the repository's quotient and
ring infrastructure, but the resulting object theorem contains only the
\ipra{} axioms.  The audit command should print the axioms of every exported
Lean theorem, exactly as in \code{IPRADerivations/Checks/Axioms.lean}.

The regular-real statements are theorem schemas.  A fixed primitive-recursive
coefficient or approximation function is introduced as a language symbol, and
its defining equations become part of \code{PRAAx}.  We never attempt to put a
Lean function variable \(\N\to\N\) inside first-order \ipra.

\subsection{Proposed module tree}

\begin{description}[style=nextline,leftmargin=0pt,labelindent=0pt,labelwidth=0pt,labelsep=0pt]
\item[\code{Algebra/Bring/Polynomial.lean}]
  Integer-difference terms for \(S_4\), the identity
  \eqref{eq:s4-sos}, strong monotonicity after importing signed order, the
  trinomial discriminant, scaling, and the differential-resolvent numerator.

\item[\code{Algebra/Bring/Rational.lean}]
  Reduced rational codes, exact order and midpoint, \(P_a(x)\), the bisection
  step, the interval invariant, \(\lceil\log_2\rceil\), and the rational
  approximation function \code{bringApprox}.

\item[\code{Algebra/Bring/RegularReal.lean}]
  The arithmetic predicate saying that the fixed approximation sequence is
  regular, compatibility of two parameter names, the Lipschitz and residual
  bounds, oddness, unboundedness, and the fifth-root asymptotic modulus.

\item[\code{Algebra/Bring/TruncatedSeries.lean}]
  Codes for vectors in \(\Q[X]/(X^{N+1})\), Cauchy product, composition,
  geometric inverse, finite residue, and the finite Lagrange theorem
  \eqref{eq:finite-lagrange}.

\item[\code{Algebra/Bring/Series.lean}]
  The Bring coefficient function, the finite identity
  \(Y_N^5+Y_N=X\bmod X^{N+1}\), the hypergeometric ratio, convergence
  majorants, and the branch-point resultant.

\item[\code{Algebra/Quintic/NormalForms.lean}]
  The depression identity, power-sum/Newton identities through degree five,
  the quadratic principal transform, generic Sylvester matrices, and the
  soundness of a Bring--Jerrard certificate.

\item[\code{Algebra/Quintic/Cyclotomic5.lean}]
  The quotient by \(1+W+W^2+W^3+W^4\), root-of-unity filters, the de Moivre
  product, the four-resolvent product, and the explicit examples.

\item[\code{Algebra/Quintic/FiniteGroups.lean}]
  Five-letter permutation codes, subgroup generation, commutators, the
  explicit proofs \([S_5,S_5]=A_5\), \([A_5,A_5]=A_5\), and the affine group
  derived series.

\item[\code{Algebra/Quintic/GaloisCertificates.lean}]
  Finite-field polynomial arithmetic, Frobenius traces, modular factorization
  certificates for the two concrete polynomials, and the finite decomposition
  group theorem needed to obtain cycle types.

\item[\code{Algebra/Quintic/RadicalTowers.lean}]
  Coded extension tables, roots of unity, Kummer kernels, solvable normal
  series, and Lagrange-resolvent reconstruction for a group contained in
  \(\operatorname{AGL}(1,5)\).
\end{description}

\subsection{First milestone: exact polynomial theorems}

This milestone uses only existing ring infrastructure and should be completed
first.  A representative object-theory definition is
\begin{lstlisting}[style=code,language={}]
-- schematic: variables are signed integer terms ZT n
abbrev S4 (x y : ZT n) : ZT n :=
  x^4 + x^3*y + x^2*y^2 + x*y^3 + y^4

def bringSOS : Fm PTm n :=
  ZT.eq (12 * S4 x y)
    (3 * (2*x^2 + x*y)^2 + (3*x*y + 2*y^2)^2 + 8*y^4)

theorem bring_sos :
  ([] : List (Fm PTm 0)) |-[PRAAx] univ n bringSOS := by
  refine closeUniv n (ZT.derive ?_)
  ipra_ring
\end{lstlisting}
The actual syntax must use the repository's \code{ZT.natLit}, multiplication
and power notation, as illustrated by
\code{Algebra/PolynomialFormulas.lean}.  The depression identity and both
cyclotomic products have the same shape: assumptions are introduced by
\code{impI}, converted into quotient-ring equalities by \code{ZT.toZ\_mem}, and
closed by \code{ipra\_ring} or \code{linear\_combination}.

The resultant identity can be proved in either of two ways:

\begin{enumerate}
\item instantiate the fixed \(9\times9\) determinant from
      \code{Core.LinearAlgebra} and normalize; or
\item prove the elimination implications in both directions using
      \(5X^4+p=0\) and \(4pX+5q=0\), together with a Bézout identity for the
      converse.
\end{enumerate}

The first route generalizes; the second produces a smaller term for the
trinomial.

\subsection{Second milestone: rational bisection}

Define total rational operations.  A denominator-zero input returns a fixed
default, and every theorem carries a positivity guard.  The bisection state is
a paired code \(\angles{L,U}\).  Its step is a primitive recursion using the
Boolean rational comparison
\[
 \chi(P_a((L+U)/2)\le0).
\]
The main induction theorem should state the three clauses of
Lemma~\ref{lem:bisection-invariant} as a single Boolean conjunction.  The
width clause is an equality after multiplying by \(2^n\), avoiding rational
exponents.

The regularity theorem for \code{bringApprox} follows from the nested-interval
invariant.  Soundness of the limit is represented arithmetically as:
for every precision \(k\), the residual of the \(k\)-th approximation is below
a displayed dyadic bound.  Equality to zero of the named real is the universal
closure of these residual bounds.  The Lipschitz theorem is first proved for
rational parameters and approximants, then transferred to the fixed regular
names by their moduli.

\subsection{Third milestone: truncated series}

A series vector can reuse the repository's pairing and bounded-sum modules.
Recommended definitions are
\begin{align*}
 \operatorname{coeff}(A,i)&=\text{the }i\text{-th decoded rational},\\
 \operatorname{mul}(A,B,n)_k
 &=\sum_{i\le k}A_iB_{k-i},\\
 \operatorname{compose}(A,B,n)
 &=\sum_{j\le n}A_jB^j\pmod{X^{n+1}}.
\end{align*}
The fixed-point proof uses the Boolean relation
\[
 \operatorname{Agree}(A,B,m)
 =\prod_{i<m}[A_i=B_i].
\]
One quantifier-free induction proves that one application of
\(A\mapsto X\Phi(A)\) increases the agreement bound by one.  A second induction
proves the finite residue substitution lemma.  These two results yield the
uniform finite Lagrange theorem.

For the Bring series, define
\begin{lstlisting}[style=code,language={}]
bringCoeff (n : Nat) : RatCode :=
  if h : n % 4 = 1 then
    let k := (n - 1) / 4
    (-1)^k * choose (5*k) k / (4*k+1)
  else 0
\end{lstlisting}
with the sign and division implemented on codes.  The object theorem is the
coefficient-vector equality, not a Lean equality of completed formal power
series.

\subsection{Fourth milestone: finite groups before Galois theory}

The finite group layer can be completed independently of field theory.  Encode
a permutation by a base-six numeral or a five-entry nested pair.  Verify
bijectivity by a five-term Boolean product.  Enumerate \(S_5\) by bounded search
through the \(5^5\) maps and filter permutations; this is primitive recursive,
although a hard-coded 120-entry table gives smaller proof terms.

The perfectness proof of \(A_5\) need not enumerate all commutators.  Formalize
the two explicit double transpositions from Lemma~\ref{lem:A5-perfect}, prove
their commutator is \((1\,2\,3)\), prove normality of the commutator subgroup,
and prove that 3-cycles generate \(A_5\).  The affine-group theorem can follow
the minimal-normal-subgroup proof of Theorem~\ref{thm:AGL-criterion}, with the
156-subgroup enumeration retained as a check theorem.

\subsection{Fifth milestone: the concrete \(S_5\) certificates}

Polynomial arithmetic modulo \(2,3,7\) is tiny.  The Frobenius remainder tables
and Bézout identities in Section~\ref{sec:abel-ruffini} can be inserted as
literal vectors and verified by \code{native\_decide} in the Lean metatheory;
a reflection lemma then turns an accepted vector computation into the
corresponding object-theory equality.

The nontrivial bridge is the factor-degree/Frobenius theorem.  A clean module
order is:

\begin{enumerate}
\item finite quotient rings and fields;
\item square-free factorization and Frobenius orbits;
\item coded finite free \(\Z\)-orders and reduction modulo a prime;
\item decomposition groups and the residue Frobenius lift;
\item cycle type on the roots.
\end{enumerate}

For the two concrete examples the splitting-field degree is at most 120, so all
objects can be bounded a priori.  The general theorem is nevertheless cleaner
than hard-coding one splitting field.

\subsection{Sixth milestone: radical towers and the generic theorem}

The finite radical-tower theorem uses basis tables and automorphism tables.
Kummer-step soundness is the equation
\((\sigma\alpha/\alpha)^n=1\).  A normal-series constructor appends the abelian
kernel at each step.  The generic \(S_5\) theorem requires a separate syntactic
ring of rational functions in five indeterminates and the constructive
symmetric-polynomial normalizer of Lemma~\ref{lem:symmetric-rational}.  It need
not quantify over arbitrary fields inside the object theory; it is a fixed
symbolic computation theorem about polynomial codes.

\subsection{Completion criteria}

A complete Lean delivery should satisfy all of the following.

\begin{enumerate}
\item \code{lake build IPRADerivations} succeeds with the new modules enabled.
\item No \code{sorry}, \code{admit}, or unproved axiom is introduced.
\item Every exported theorem has the closed \code{Prf PRAAx} type or is clearly
      marked as a metatheoretic helper.
\item \code{Checks/Axioms.lean} reports only the same Lean metatheory axioms
      already present in the repository, and no new object-theory axiom.
\item The Python certificate JSON is regenerated and its hashes match the
      archive manifest.
\item The formal statement audit points to the exact Lean theorem name for
      every item upgraded from \(\Outline\) to \(\Exact\) or
      \(\Effective\).
\end{enumerate}

The ordering above deliberately yields useful checked theorems early: the
polynomial identities, real bisection, series coefficients and finite groups do
not have to wait for the larger Galois infrastructure.
